\scanpage{330}\mspage{246}
% l. 1–4
\unsure{De ceci|D'ici}, le th. 4.10.1 résulte de façon essentiellement \unsure{verbale}, tout comme dans le cas particulier 4.9.3 précédent (cf page 223). \deleted{\unsure{Pour}}
% slanted
\margin{un peu court ! Il faudrait expliciter l'argument ou en faire \unsure{un}, il y en a pour \unsure{deux} \unsure{ou} trois pages\unsure{!}}

% l. 5–11
Pour établir le corollaire, comme la catégorie-but de (4.10.2.2) est \unsure{grde|gde} devant $\unsure{\pi|\kappa}$, comme produit de telles catégories (cf 4.9.6.6 page 228), \inserted{et le foncteur est fidèle,} \inserted{il suffit} en vertu de 4.9.6.3, $2^\circ$, de vérifier les conditions $\mathrm F1$ et $\mathrm F3_{\pi}$ \deleted{\unsure{pour}} (page 225) pour le foncteur envisagé.

% l. 12–15
On va interpréter ce foncteur (4.10.2.2) en termes d'une catégorie cofibrée
\[
 \mathscr D_{\unsure{\varphi|\phi}}\longrightarrow J, \tag*{(4.10.5)}
\]
dont les fibres sont les $\mathscr D_j/\varphi(j)$.
% l. 16–19
On note que si $\alpha:j\unsure{\mathbin{-}|\longrightarrow}j'$ est une flèche de $J$, alors on a
\[
 \left\{\begin{aligned}
 \alpha_* &: C_j\unsure{\mathbin{-}|\longrightarrow} C_{\unsure{j'|j}},\\
 &\text{et }\alpha_*(\varphi(j))\longrightarrow\varphi(j')
 \end{aligned}\right.
\]
% archive sheet number: 329
